\documentclass[11pt]{article}
\usepackage[a4paper,margin=18mm]{geometry}
\usepackage{fontspec}
\setmainfont{Latin Modern Roman}
\usepackage{amsmath,amssymb,amsthm,amscd,graphicx,color}
\usepackage{xurl}
\usepackage[unicode,hidelinks]{hyperref}
\allowdisplaybreaks
\setlength\emergencystretch{3em}
\numberwithin{equation}{section}

\newcommand{\Hf}{\mathop\mathrm{Hf}\nolimits}
\newcommand{\Pf}{\mathop\mathrm{Pf}\nolimits}
\newcommand{\Res}{\mathop\mathrm{Res}\nolimits}
\newcommand{\Ber}{\mathop\mathrm{Ber}\nolimits}
\newcommand{\Z}{\mathbb{Z}}
\newcommand{\INZ}[2]{\in\frac#1#2 +\mathbb{Z}}
\newcommand{\C}{\textsf{C}_}

\newcommand\g{\gamma}
\newcommand{\G}{\Gamma}
\newcommand{\bt}{{\bf {t}}}

\renewcommand{\l}{\langle}
\renewcommand{\r}{\rangle}
\newcommand{\ve}{\varepsilon}



\begin{document}
\begin{center}\Large The negative-level CKP free-boson Fock module has a graded Heisenberg-plus-neutral-fermion polynomial realization with the explicit dressed-field vertex operator\end{center}
\noindent\textbf{Stable ID:} 2141\par
\noindent\textbf{Authors:} Johan W. van de Leur; Alexander Yu. Orlov; Takahiro Shiota\par
\noindent\textbf{Original work:} CKP Hierarchy, Bosonic Tau Function and Bosonization Formulae\par
\noindent\textbf{Version:} Official SIGMA 8 (2012) article 036, DOI 10.3842/SIGMA.2012.036; journal PDF and native TeX acquired 2026-10-01, exact bytes pinned by SHA256\par
\noindent\textbf{Selected task scope:} Full unlabelled construction in Sections2.1-2.3: algebraic free-boson Fock space F with half-integer modes phi\_j, commutation (2.1), vacuum conditions (2.2), ordered monomial basis (2.3) and half-integer grading. Odd quadratic currents J\_n satisfy (2.7); the dressed field theta(z)=\allowbreak{}V\_-(z)\textasciicircum{}(-1)phi(z)V\_+\allowbreak{}(z)\textasciicircum{}(-1) has neutral-fermion modes commuting with the currents and supercommutation (2.18). The graded isomorphism sigma identifies F with the polynomial superalgebra C[t\_(2j-1),t\_((2j-1)/2):j>=\allowbreak{}1], with commuting integer-odd times and Grassmann positive-half-integer times, mode actions (2.19), and the formal boson vertex (2.20). All field identities are coefficientwise; no analytic completion, ordinary positive-level KP/CKP correspondence or hierarchy equality is claimed.\par
\noindent\textbf{Difficulty:} H2: The correspondence requires constructing a dressed neutral fermion commuting with the odd Heisenberg currents and evaluating its anticommutator by a complete formal-delta-function contraction calculation. Together with the graded partition-character factorization, this identifies the entire free-boson Fock space with mixed commuting and Grassmann variables and yields the original boson vertex field, beyond ordinary qualification Lie representation techniques.\par
\noindent\textbf{Editorial boundary note (not author text):} The complete original unlabelled Section 2.3 Bose-to-Heisenberg-plus-neutral-fermion correspondence is retained with its graded sigma isomorphism, exact mode actions (2.19) and dressed boson vertex formula (2.20). All Sections 2.1-2.3 construction, negative-level odd quadratic currents, character factorization, dressing, commuting-current proof, three reordering identities, full formal-delta anticommutator and supermode relations are supplied as one main construction unit. The Fock space and polynomial realization are algebraic, with the vertex field interpreted coefficientwise formally. Standard Kac formal-delta identities and bounded mode/graded-character extraction remain the author prerequisites and checks. This is the negative-level CKP construction, not a claim identifying it with the usual Lax-reduction hierarchy. No tau/bilinear, Pfaffian/Hafnian or polynomial-family result is separately selected. The later original notation formula for Gamma(t), H(t) and chi(t\_odd) is supplied solely to resolve the author warning that H(y) in the anticommutator is a different field, not a missing new proof obligation.\par
\noindent\textbf{Source:} \url{https://sigma-journal.com/2012/036/}\quad DOI: \url{https://doi.org/10.3842/SIGMA.2012.036}\par
\noindent\textbf{License and attribution:} Copyright retained by the named authors. Published by SIGMA under CC BY 4.0: \url{https://creativecommons.org/licenses/by/4.0/}. Source: \url{https://sigma-journal.com/about.html}. Adaptation consists of standalone excerpt formatting, cover and numbering restoration; original author language and mathematical text are retained.\par
\medskip\hrule\medskip\noindent\textbf{Author text begins below.}\par
\setcounter{section}{1}
\setcounter{subsection}{0}
\setcounter{subsubsection}{0}
\setcounter{equation}{0}
% Original source lines 116--469
\section{CKP bosonic tau function\label{CKPbosonic}}

In this section we follow a suggestion in~\cite{DJKMVI} and describe a CKP hierarchy of PDEs starting from  a collection of free bosons.
This imitates their approach in the BKP case, where one starts with neutral fermions.
However, this hierarchy, although related to Lie algebra $c_\infty$,
dif\/fers from the usual CKP hierarchy, for which one takes
a reduction of the KP hierarchy by assuming that the Lax operator
satisf\/ies $L^*=-L$; such Lax operators come from certain KP tau functions,
which are f\/ixed by some involution and where one puts the even times to zero,
see e.g., \cite{DJKMVI} or \cite{AL} for more details.
The hierarchy described in this paper is dif\/ferent and is not related to the
usual CKP which comes from a reduction of KP. In the latter case
one has a realization of $c_\infty$, for which the level is positive.
Our construction realizes  $c_\infty$ with a {\it negative\/} level.

\subsection[Bose-Fermi correspondence in the CKP case]{Bose--Fermi correspondence in the CKP case}\label{Bose-Fermi}

We follow a suggestion of Date, Jimbo, Kashiwara and Miwa in their paper \cite{DJKMVI}
and introduce free bosons, but for convenience of notation we
shift the index by $\frac 12$. So
$\phi_i$ with $i\in\frac 12+\mathbb{Z}$ satisfy commutation relations:
\begin{gather}
\label{S1}
\phi_{i}\phi_j-\phi_j\phi_i=(-)^{j-\frac{1}{2}}\delta_{i,-j}.
\end{gather}
The Fock space $F$, respectively $F^*$, is def\/ined
by
\begin{gather}
\label{S2} \phi_j\, |0\rangle=0\quad\mbox{if}\ \ j<0,\qquad\mbox{resp.}\qquad
\langle 0|\, \phi_j=0\quad\mbox{if}\ \ j>0,
\end{gather}
so that $F$ has as basis the vectors
 \begin{gather}\label{S2BASIS}
(\phi_{j_1})^{m_1}(\phi_{j_{2}})^{m_{2}}
\cdots(\phi_{j_{n-1}})^{m_{n-1}}(\phi_{j_n})^{m_n}|0\rangle
 \end{gather}
with $j_1>j_2>\cdots> j_{n-1}>j_n>0$ and $m_i$ positive integers.
Def\/ining
\[
\deg |0\rangle =0 ,\qquad\deg\phi_j=j ,
\]
we have a direct sum decomposition of $F$:
\[
F=\bigoplus_{k\in\frac12\mathbb Z} F_k\qquad\mbox{with}\quad
F_k=\{f\in F\mid \deg f=k\} .
\]
It is straightforward to check that the dimension of $F_k$ is
given by the partition of $k$ into
 positive elements of $\frac 12+\mathbb{Z}$.
Def\/ine the formal character as
\[
\dim_q F=\sum_{k\in\frac12\mathbb Z} \dim F_k\, q^k .
\]
Then
\begin{gather}\label{q1}
\dim_q F=\prod_{0<k\INZ12}\frac 1{1-q^k} .
\end{gather}
Writing
\[
\phi(z)=\sum_{j\INZ12}\phi_jz^{j-\frac 12} ,
\]
we denote
\begin{gather}
\label{S4} H(z):=\sum_{n\in
1+2\mathbb{Z}}J_nz^{-n-1}:=-\frac 12\,{:}\phi(-z)\phi(z){:} ,
\end{gather}
where the normal ordering is def\/ined by
\begin{gather}\label{normal ordering}
{:}\phi_i\phi_j{:}=
\begin{cases}
\phi_i\phi_j&\mbox{if }i\ge j ,\\[2pt]
\phi_j\phi_i&\mbox{if }j>i .
\end{cases}
\end{gather}
In other words  $J_n=0$ for $n$ even and
\[ % \begin{gather} \label{S}
J_n=\frac{1}{2}\sum_{j\in
\frac{1}{2}+\mathbb{Z}}(-)^{j+\frac{1}{2}}\phi_j\phi_{-j-n}\qquad
\mbox{for\,\ $n$\,\ odd;}
\] % \end{gather}
one has the following familiar commutation relations
\begin{gather}
\label{S5} [J_n, J_m]=-\frac{n}{2}\delta_{m,-n} .
\end{gather}
The elements ${:}\phi_i\phi_j{:}$ form a representation of the
Lie algebra $c_\infty$, see e.g.~\cite{Kac-Leur}. However,
we want to stress that its level (the value of its central element) is negative.
Note also that in the commutation relations~(\ref{S5}) we
have the factor $-\frac{n}{2}$ instead of the usual~$\frac{n}{2}$.

It is clear that
\[
J_n|0\rangle=\langle 0|J_{-n}=0\qquad\mbox{for}\ \ n>0 .
\]
By a similar argument as before,
again since these are bosons, we can apply an element $J_{-n}$
inf\/initely many times to $|0\rangle$. Since the degree of $J_n$ is $-n$
 we obtain that the action of this Heisenberg algebra on the vacuum
vector produces in the $\dim_q F$ the partition function of partitions in only odd numbers:
\[
\prod_{0< k\in 1+2\mathbb{Z}}\frac{1}{1-q^k}.
\]
Now we calculate, using \eqref{q1},
\begin{gather}
\left(\prod_{0< k\in
1+2\mathbb{Z}}\frac{1}{1-q^k}\right)^{-1}\dim_q F=
\prod_{0< k\in\frac{1}{2}+\mathbb{Z}}\frac{{1-q^{2k}}}{1-q^k}
=\prod_{0< k\in\frac{1}{2}+\mathbb{Z}}\big(1+q^k\big).\label{q3}
\end{gather}
 This part should be explained by something else and we expect it to be
fermions, at least anticommuting variables. The factor~$1+q^k$
is related to a fermion of degree~$k$. This is how we
get these elements and calculate their commutation relations.

We f\/irst calculate
\begin{gather}
  [J_n, \phi(z)] =\frac 12\sum_{j,k\INZ12}(-)^{j+\frac 12}
[\phi_j\phi_{-j-n}, \phi_k]z^{k-\frac 12}\nonumber\\
\hphantom{[J_n, \phi(z)]}{} =\frac{1}{2}\sum_{j,k\INZ12}(-)^{j+\frac 12} \left(
[\phi_j, \phi_k]\phi_{-j-n}+ \phi_j[\phi_{-j-n}, \phi_k]
\right) z^{k-\frac12}\nonumber\\
\hphantom{[J_n, \phi(z)]}{}
=\frac 12\sum_{j,k\INZ12}(-)^{j+\frac 12}
(-)^{k-\frac 12}\left(\delta_{j,-k}\phi_{-j-n}+\delta_{j+n,k}\phi_j\right)
z^{k-\frac 12}\nonumber\\
\hphantom{[J_n, \phi(z)]}{}
=\frac12\sum_{k\INZ12} 2\phi_{k-n}z^{k-\frac 12}
=z^n\phi(z) .\label{S6}
\end{gather}
Now, using (\ref{S5}) we see that
\begin{gather}\label{S5bis}
\big[J_n,e^{\frac2m J_mz^{-m}}\big]=\delta_{n,-m}z^n e^{\frac2m J_mz^{-m}}.
\end{gather}
Hence setting
\begin{gather}\label{theta}
 \theta(z):=V_-(z)^{-1}\phi(z)V_+(z)^{-1},
\end{gather}
where
\begin{gather}
\label{exp}
V_\pm(z)
=\exp\sum_{\pm k>0,\,\mathrm{odd}}\frac{2}{k}J_k z^{-k} ,
\end{gather}
we have, from
\eqref{S6} and \eqref{S5bis},
\begin{gather}\label{comm-Htheta}
[J_n, \theta(z)]=0 .
\end{gather}

\subsection[Commutation relations of the $\theta(z)$'s]{Commutation relations of the $\boldsymbol{\theta(z)}$'s}

We now want to calculate the commutation relations of these
$\theta(z)$'s given in (\ref{theta}).  For this we f\/irst rewrite the commutation relations
(\ref{S1}) as follows:
\[
\phi(z)\phi(y)-\phi(y)\phi(z)=\delta(z-(-y))\,,
\]
where
\[
\delta(z-y)=z^{-1}\sum_{k\in\mathbb{Z}}
\left(\frac{z}{y}\right)^k.
\]
Note also that
\[
\phi(-z)\phi(y)={:}\phi(-z)\phi(y){:}-\frac{1}{z}\frac{1}{1-\frac{y}{z}}  .
\]
We f\/irst show the following identities
\begin{gather}
V_+(-z)^{-1}V_-(y)^{-1} =
\frac{1-\frac{y}{z}}{1+\frac{y}{z}}V_-(y)^{-1}V_+(-z)^{-1},\nonumber\\
\phi(-z)V_-(y)^{-1} =
\frac{1+\frac{y}{z}}{1-\frac{y}{z}}V_-(y)^{-1}\phi(-z),\nonumber\\
V_+(-z)^{-1}\phi(y) =
\frac{1+\frac{y}{z}}{1-\frac{y}{z}}\phi(y)V_+(-z)^{-1}.\label{ident}
\end{gather}
Introduce
\[
V(t)=\exp\sum_{k>0,\,\mathrm{odd}}t_k J_k .
\]
Then
\begin{gather}
\label{VT1}
V_+(z)=V\left(\frac21 z^{-1},\frac23 z^{-3},\frac25 z^{-5},\ldots\right) .
\end{gather}
The f\/irst equation of (\ref{ident}) is obtained in the following
way. First using \eqref{S5} one has
\begin{gather*}
V(t)V_-(y) =\exp \left[  \sum_{k>0,\,\mathrm{odd}} t_kJ_k,
- \sum_{\ell>0,\,\mathrm{odd}} \frac{2}{\ell}y^{\ell}J_{-\ell}\right]V_-(y)V(t)\\
\hphantom{V(t)V_-(y)}{}  =\exp \left(\sum_{k>0,\,\mathrm{odd}}t_{k}y^k
\right)
V_-(y)V(t)  .
\end{gather*}
Combining this with (\ref{VT1}) one obtains
\begin{gather*}
V_+(-z)^{-1}V_-(y)^{-1} =\exp \left(-\sum_{k>0,\,\mathrm{odd}} \frac{2}{k}\left(
\frac{y}{z}\right)^k\right)
V_-(y)^{-1}V_+(-z)^{-1}\\
\hphantom{V_+(-z)^{-1}V_-(y)^{-1}}{}
 =\frac{1-\frac{y}{z}}{1+\frac{y}{z}}V_-(y)^{-1}V_+(-z)^{-1}.
\end{gather*}
Using $\phi(-z)J_n=(J_n-(-z)^n)\phi(-z)$, see~\eqref{S6},
we obtain the second relation in \eqref{ident} as follows:
\begin{gather*}
\phi(-z)V_-(y)^{-1} =
\phi(-z)\exp\left(\sum_{k>0,\,\rm odd}\frac2kJ_{-k}y^k\right)
 =
\exp\left(\sum_{k>0,\,\rm odd}\frac2k(J_{-k}+z^{-k})y^k\right)\phi(-z)
\\
 \phantom{\phi(-z)V_-(y)^{-1}}{}
 =V_-(y)^{-1}\exp \left( \sum_{k>0,\,\mathrm{odd}}
\frac{2}{k}\left( \frac{y}{z}\right)^{k}\right)\phi(-z)
 =\frac{1+\frac{y}{z}}{1-\frac{y}{z}}V_-(y)^{-1}\phi(-z).
\end{gather*}
The third formula is proved in a similar way.

We will also use the
following identities which can be found in V.~Kac's book~\cite{KacVAfB}:
\begin{gather*}
(z-y)\partial_y\delta(z-y) =\delta(z-y), \\
(z-y)^{k+1}\partial_y^{k}\delta(z-y) =0, \\
\delta(z-y)a(z) =\delta(z-y)a(y), \\
\partial_y\delta(z-y)a(z) =\partial_y\delta(z-y)\left(a(y)+(z-y)\partial_ya(y)\right).
\end{gather*}

We now calculate
\begin{gather*}
\theta(-z)\theta(y) =V_-(-z)^{-1}\phi(-z)V_+(-z)^{-1}V_-(y)^{-1}\phi(y)V_+(y)^{-1}\\
\hphantom{\theta(-z)\theta(y)}{} =\frac{1-\frac{y}{z}}{1+\frac{y}{z}}
V_-(-z)^{-1}\phi(-z)V_-(y)^{-1}V_+(-z)^{-1}\phi(y)V_+(y)^{-1}\\
\hphantom{\theta(-z)\theta(y)}{}=V_-(-z)^{-1}V_-(y)^{-1}\phi(-z)V_+(-z)^{-1}\phi(y)V_+(y)^{-1}\\
\hphantom{\theta(-z)\theta(y)}{} =\frac{1+\frac{y}{z}}{1-\frac{y}{z}}
V_-(-z)^{-1}V_-(y)^{-1}\phi(-z)\phi(y)V_+(-z)^{-1}V_+(y)^{-1}\\
\hphantom{\theta(-z)\theta(y)}{} =\frac{1+\frac{y}{z}}{1-\frac{y}{z}}
V_-(-z)^{-1}V_-(y)^{-1}\left({:}\phi(-z)\phi(y){:}-\frac{1}{z}\frac{1}{1-\frac{y}{z}}\right)
V_+(-z)^{-1}V_+(y)^{-1}.
\end{gather*}
Now replacing $z$ and $y$ by $-y$ and $-z$ respectively, gives
\[
\theta(y)\theta(-z)=\frac{1+\frac{z}{y}}{1-\frac{z}{y}}
V_-(-z)^{-1}V_-(y)^{-1}\left({:}\phi(-z)\phi(y){:}+\frac{1}{y}\frac{1}{1-\frac{z}{y}}\right)
V_+(-z)^{-1}V_+(y)^{-1}
\]
and thus
 \begin{gather}
 \theta(-z)\theta(y)+\theta(y)\theta(-z)
=
2z\delta(z-y)V_-(-z)^{-1}V_-(y)^{-1}{:}\phi(-z)\phi(y){:}
V_+(-z)^{-1}V_+(y)^{-1}\nonumber\\
\qquad\quad{} + \left( \frac 1y \frac{1+\frac{z}{y}}{(1-\frac{z}{y})^2}-\frac
1z \frac{1+\frac{y}{z}}{(1-\frac{y}{z})^2}\right)
V_-(-z)^{-1}V_-(y)^{-1} V_+(-z)^{-1}V_+(y)^{-1}\nonumber\\
\qquad {}
= 4yH(y)\delta(z-y)-\partial_y\delta(z-y)(y+z)V_-(-z)^{-1}V_-(y)^{-1}
V_+(-z)^{-1}V_+(y)^{-1}\nonumber\\
\qquad {}
= 4yH(y)\delta(z-y)-2y\partial_y\delta(z-y)
-(z-y)\partial_y\delta(z-y) \nonumber\\
\qquad\quad{}
 \times\left(1+2y\partial_y\left(V_-(-y)^{-1}\right)V_-(y)^{-1}
+2y\partial_y\left(V_+(-y)^{-1}\right)V_+(y)^{-1}\right)\nonumber\\
\qquad {}
= 4yH(y)\delta(z-y)-2y\partial_y\delta(z-y)-
\delta(z-y)\left(1+4yH(y)\right)\nonumber\\
\qquad {}
= {-2}y\partial_y\delta(z-y)-\delta(z-y)=-D_y\delta(z-y) ,\label{thetatheta-equality}
\end{gather}
where $H(y)$ is as in \eqref{S4} (not the one in \eqref{Gamma-bt2}),
and $D_y=y\partial_y+\partial_y y$ is the Euler operator.



Now write
\begin{gather}\label{theta(z)}
\theta(z)=2\sum_{i\in\frac{1}{2}+\mathbb{Z}}
J_iz^{-i-\frac{1}{2}} .
\end{gather}
Note that there is no conf\/lict with the $J$'s def\/ined in (\ref{S4}), since here the $J$'s
have indices in $\frac{1}{2}+\mathbb{Z}$.
It is clear that the above commutation
relation \eqref{thetatheta-equality} in modes gives
\[
J_jJ_k+J_kJ_j=(-)^{j-\frac{1}{2}}\frac{j}2 \delta_{j,-k},\qquad j,\,k\in \frac{1}{2}+\mathbb{Z}
\]
(compare with (\ref{S5})). From \eqref{comm-Htheta} we also have
\[
[J_n, J_m]=0 ,\qquad n\in 1+2\mathbb{Z},\quad  m\in \frac{1}{2}+\mathbb{Z} .
\]
Thus we can combine the (anti)commutation relations of all $J$'s as follows:
 \begin{gather}\label{supercurrent}
[J_i, J_j]_s = \frac{j}{2}(-1)^{[j-\frac 12]}\delta_{i,-j} ,
 \end{gather}
where the notation $[\ {,}\ ]_s$ serves for the supercommutator while
$[i]$ denotes the integer part of a~real number~$i$. As we see,
$\deg J_i=-i$
 and that
\[
J_k|0\rangle=\langle 0|J_{-k}=0,\qquad
J_{-k}|0\rangle\ne 0\ne\langle 0|J_{k} \qquad\mbox{for}\ \ k>0 .
\]

\subsection{Even and odd times} \label{Even and odd times}
Since the elements $J_{-k}$ with $k\in \frac{1}{2}+\mathbb{Z}$
anticommute among themselves, they can only
appear once in
\[
J_{-k_n}J_{-k_{n-1}}\cdots J_{-k_3}J_{-k_2}J_{-k_1}|0\rangle.
\]
Such a $J_{-k}$ explains the factor $1+q^k$ in the
$q$-dimension formula \eqref{q3}. One can identify the~$J_n$'s, for $n<0$ with even and odd times, i.e., with  commuting
variables $t_j$, $0<j\in 1+2\mathbb{Z}$, and Grassmann variables $t_{\frac{j}2}$, $0<j\in
1+2\mathbb{Z}$, and identify Fock space $F$ with the space
\[
\mathbb{C}\big[t_{2j-1},t_{\frac{2j-1}2};\, j=1,2,\ldots \big]
\]
(or some completion of it, since we take exponentials),
where one has
\[
t_it_j-(-)^{4ij}t_jt_i=0 ,
\]
in particular, $t_j^2=0$ for $j\in \frac12+\mathbb{Z}$\,.
We will write ${\bf t}=\bigl(t_1,t_{\frac 12};t_3,t_{\frac 32};t_5,t_{\frac
52};\dots \bigr)$ and use $t=\left(t_1,t_3,t_5,\dots \right)$ and
$t_\mathrm{odd}=\bigl(t_{\frac 12},t_{\frac 32};,t_{\frac
52},\dots \bigr)$.


{\it Let $\sigma$ be this isomorphism,
sending $F$ to $\mathbb{C}[t_{2j-1},t_{\frac{2j-1}2}; j=1,2,\ldots]$.}
Then
\begin{gather}\label{theta-xi}
\sigma J_{-j}\sigma^{-1}=(-)^{[\frac{1}{2}-j]}\frac{j}2t_j\qquad\mbox{and}\qquad
\sigma J_j\sigma^{-1}=\frac{\partial}{\partial t_j}\,,\qquad
j>0\,,
\end{gather}
give the f\/ield exactly in commuting and anticommuting variables
$t_k$.


Now using the free boson-(boson+fermion) correspondence, i.e.,
using the vertex operator expressions for the f\/ields
\begin{gather}
\sigma\phi(z)\sigma^{-1} =
\exp\left( \sum_{\scriptstyle0<k\in\mathbb{Z}\atop\scriptstyle k:\,\mathrm{odd}}t_kz^k \right)
\exp\left( \sum_{\scriptstyle0<k\in\mathbb{Z}\atop\scriptstyle k:\,\mathrm{odd}}
\frac{2}{k}\frac{\partial}{\partial t_k}z^{-k} \right)\nonumber\\
\hphantom{\sigma\phi(z)\sigma^{-1} =}{}\times
\sum_{0<j\in\mathbb{Z}}\left( (2j-1)t_{\frac{2j-1}2}(-z)^{j-1}+2\frac{\partial}{\partial
t_{\frac{2j-1}2}}z^{-j} \right)\label{phixi}
\end{gather}
(where we used (\ref{exp}), (\ref{theta}), (\ref{theta(z)}) and (\ref{theta-xi})),
in the following subsections
we shall express the bilinear identity as a hierarchy of dif\/ferential equations.
A similar expression  for~(\ref{phixi}) was also found in~\cite{Kac-Leur}.
\setcounter{section}{2}
\setcounter{subsection}{5}
\setcounter{subsubsection}{0}
\setcounter{equation}{27}
% Original source lines 622--626
\begin{gather}\label{Gamma-bt2}
\Gamma({\bf t})=e^{H(t)}e^{\chi(t_\mathrm{odd})} ,\qquad  \mbox{where}\quad
H(t):=\sum_{0< i\in 1+2\Z}  t_i J_{i} ,\qquad
\chi(t_\mathrm{odd}):=\sum_{0< i\INZ12}  t_i J_{i} .
\end{gather}
\begin{thebibliography}{99}
\bibitem[1]{AL}
Aratyn H., van~de Leur J.W., The {CKP} hierarchy and the {WDVV} prepotential, in
  Bilinear Integrable Systems: from Classical to Quantum, Continuous to
  Discrete, \href{http://dx.doi.org/10.1007/978-1-4020-3503-6_1}{\textit{NATO Sci. Ser.~II Math. Phys. Chem.}}, Vol.~201, Springer,
  Dordrecht, 2006, 1--11, \href{http://arxiv.org/abs/nlin.SI/0302004}{nlin.SI/0302004}.

\bibitem[3]{DJKMVI}
Date E., Jimbo M., Kashiwara M., Miwa T., Transformation groups for soliton
  equations. {VI}.~{KP} hierarchies of orthogonal and symplectic type,
  \href{http://dx.doi.org/10.1143/JPSJ.50.3813}{\textit{J.~Phys. Soc. Japan}} \textbf{50} (1981), 3813--3818.

\bibitem[8]{KacVAfB}
Kac V., Vertex algebras for beginners, \textit{University Lecture Series},
  Vol.~10, 2nd ed., American Mathematical Society, Providence, RI, 1998.

\bibitem[9]{Kac-Leur}
Kac V.G., van~de Leur J.W., Super boson-fermion correspondence of type~$B$, in
  Inf\/inite-Dimensional {L}ie Algebras and Groups ({L}uminy-{M}arseille, 1988),
  \textit{Adv. Ser. Math. Phys.}, Vol.~7, World Sci. Publ., Teaneck, NJ, 1989,
  369--406.

\end{thebibliography}
\end{document}
